\documentclass[12pt]{article}

\usepackage{amsmath}
\usepackage{german}
\usepackage{graphicx}
\DeclareMathOperator\sgn{sgn}

%\graphicspath{{/home/ritz/Teaching/fcn/fig/}}

%----------------- layout ----------------------------------------

\oddsidemargin=0in
\textwidth=6in
\topmargin=-0.7in
\textheight=9.5in
\baselineskip=16pt
\font\mittel=cmr17
\font\roman=cmr12
\font\klein=cmr9
\font\mgross=cmr10
\newfont{\CAP}{cmcsc10 scaled 900}
\pagestyle{empty}
\parindent=0mm

%----------------- macros ----------------------------------------
\newcommand{\conc}[1]{[{\rm C}]_{\text{#1}}}
\newcommand{\units}[1]{\text{#1}}
\newcommand{\K}{\text{K}}
\newcommand{\Na}{\text{Na}}
\newcommand{\Cl}{\text{Cl}}
\newcommand{\Leak}{\text{leak}}
\newcommand{\water}{\text{H}_2\text{O}}
\newcommand{\tderiv}[1]{\frac{\mathrm{d}#1}{\mathrm{d}t}}
%-----------------------------------------------------------------


\begin{document}

\parbox{2cm}{
\includegraphics[width=1.8cm]{../bccnlogo.pdf}
}
\parbox{11cm}{
\begin{center}
\large HUMBOLDT-UNIVERSIT"AT \hskip 0.1 cm ZU \hskip 0.1 cm BERLIN
\vskip 0.1 true cm
\mgross BERNSTEIN CENTER FOR\\ COMPUTATIONAL NEUROSCIENCE
\end{center}
}
\parbox{2cm}
{
\hfill
\includegraphics[width=1.8cm]{../hublogo.pdf}
}

\vskip 0.6 true cm

\leftline{\CAP Humboldt-Universit"at zu Berlin 
\hfill Phone: 030/2093-9110}
\leftline{\CAP Philippstr. 13 House 6
             \hfill Fax: 030/2093-6771}
\leftline{\CAP \hfill webpage: http://www.bccn-berlin.de/}


\vskip 0.8 true cm
\centerline{\bf Models of Neural Systems, WS 2009/10}
\centerline{\bf Project 1: Phase oscillators} 
\centerline{Project presentation and report submission: February,
8th, 2010}

\vskip 0.8 true cm

{\bf Background}

Neural systems quite often reveal periodic dynamics e.g. regular firing
neurons or oscillatory populations of neurons. The details of
the dynamics of such models tend to be quite complicated which makes
it hard to analyse them. Nevertheless, it can be shown that  the
dynamics  in the limit of small perturbations can be reduced to so-called phase
oscillators in which only the phase of the oscillation is modelled and
the amplitude is assumed to be constant. The very simplest models has
the following form:


\begin{equation}
    \tderiv{\phi}(t)=\omega
    \label{<++>}
\end{equation}

Here you will study synchronisation properties of such phase
oscillators.
\vskip 0.8 true cm

{\bf Problems}

\begin{enumerate}
    \item Implement a single oscillator perturbed by an external
        stimulus $I(t)$. Lets assume that the effect of the stimulus
         on the dynamics is phase-dependent,
        so that we can write:
        \begin{equation}
            \tderiv{\phi_i} = \omega_i + I(t) \sin{\phi_i}
            \label{<++>}
        \end{equation}
        where the stimulus is a short pulse starting at time
        $t_\mathrm{onset}$ with duration $T$ and intensity of $I_0$:
        \begin{equation}
          I(t)=
          \begin{cases}
              I_0, & \text{if } t_\mathrm{onset}+T>t>t_{\mathrm{onset}} \\
              0, & \text{otherwise} 
          \end{cases}
            \label{eq:stimulus}
        \end{equation}
        Plot phase-response curve i.e. the post-stimulus phase, as a
        function of the pre-stimulus phase for different $I$ and
        $t_\mathrm{onset}$.
    \item Study the effect of the stimulus on a population of
        $N=100$ identical phase oscillators (no coupling) with random initial
        conditions. Plot the phases as a color image. Calculate and
        plot the order parameter:
        \begin{equation}
            Z = \left|{\sum_i^N \exp{(j\phi_i)}}\right|,
            \label{eq:order_parameter}
        \end{equation}
        where $j=\sqrt{-1}$ is an imaginary unit. What does the
        order parameter quantify?
    \item Introduce white noise to the model:
        \begin{equation}
            \tderiv{\phi_i} = \omega + I(t)\sin{\phi_i} + \sigma \eta(t),
            \label{eq:phase_with_noise}
        \end{equation}
        where $\eta(t)$ is a normally distributed random variable of
        unit standard deviation and zero mean. Study the effects of
        the noise on the dynamics and synchronisation properties.
    \item \textit{Kuramoto mode} Simulate and analyse a population of
        globally coupled phase oscillators without noise. The dynamics of each
        oscillator is described by:
        \begin{equation}
            \tderiv{\phi_i} = \omega_i + K/N\sum_{j=1}^N\sin(\phi_j-\phi_i),
            \label{eq:kuramoto}
        \end{equation}
        where $\omega_i$ is an uniformly distributed random variable.

        Under which conditions can the oscillators synchronize? Plot
        the order parameter $Z$ as a function of the bifurcation parameter $K$.
    \item It has been shown that oscillators can also form distinct
        clusters of synchronised activity. It can be achieved in a
        model when a non-trivial coupling is chosen. Read the paper
        Tass, Stochastic phase resetting of stimulus-locked responses of two coupled oscillators: transient response clustering, synchronization, and desynchronization. Chaos vol. 13 (1) pp. 364-376
        and try to reproduce its results.
\end{enumerate}

Richard Kempter \hfill Phone: 2093-8925 \hfill
Email: r.kempter(at)biologie.hu-berlin.de \\
Robert Schmidt \hfill Phone: 2093-8926 \hfill
Email: r.schmidt@biologie.hu-berlin.de \\
Bartosz Telenczuk \hfill Phone: 2093-8838 \hfill
Email: b.telenczuk@biologie.hu-berlin.de \\

\end{document}
